\clearpage\fancyhead[L]{\small COMET.}\hypertarget{native-378-lower}{}\label{native-378-lower}
\begin{center}\textbf{COMET.}\par\vspace{4pt}\hrule\vspace{4pt}\fontsize{10}{12}\selectfont Göttingische gelehrte Anzeigen. Stück 19. S. 185\;..\;189. 1814. Januar 31.\par\vspace{4pt}\hrule\vspace{4pt}\end{center}

\par\indent Aus einem Schreiben des Hrn. Prof. \textsc{Bessel} in Königsberg an den Prof. \textsc{Gauss} vom 30. December v. J. heben wir hier einiges aus, was die astronomischen Leser unsrer Blätter gern früh zu ihrer Kenntniss gebracht sehen werden. Es ist bereits aus verschiedenen Nachrichten bekannt, dass die auch in unsern Anzeigen mehrere Male angedeutete Aussicht (man vergl. 1811, St. 130 und 1812, St. 67 [S. 327 und 349 d. B.]), der grosse Comet von 1811 könne wohl noch einmal im Sommer des Jahrs mit vorzüglichen Instrumenten und unter günstigen Umständen wieder gesehen und beobachtet werden, nicht grundlos gewesen ist. Allein nicht den, wenn gleich sorgfältigen, Bemühungen der Astronomen auf den Europäischen Sternwarten ist dies Glück zu Theil geworden, eine Folge theils der in jener Jahrszeit zu starken nächtlichen Dämmerung, theils des weniger günstigen Climas, womit die meisten dieser Astronomen zu kämpfen haben. Vom Caucasus her, aus \emph{Neu-Tscherkask} erhalten die Astronomen die in ihrer Art einzigen Beobachtungen des merkwürdigen Cometen, fast ein Jahr nach seinem Durchgange durch die Sonnennähe von Herrn von \textsc{Wisniwsky} angestellt, welcher schon seit mehrern Jahren jene entlegenen Gegenden bereist. Hr. Prof. \textsc{Bessel}, welcher sich bekanntlich um die Theorie dieses Cometen schon in mehr als einer Hinsicht verdient gemacht, und eine elliptische Bahn für denselben berechnet hatte, er\SourceLeaf{379}hielt diese wichtigen Beobachtungen von Petersburg aus mitgetheilt, und wird dieselben künftig zur feinsten Ausfeilung der Bahn benutzen. Folgende vorläufige Resultate verdienen indess schon jetzt bekannt gemacht zu werden.

\Needspace{14\baselineskip}\par\indent Herr von \textsc{Wisniwsky} beobachtete den Cometen vom 8. bis 27. August, er erhielt in dieser Zeit 29 Durchgänge durch das Kreismikrometer, dessen Feld er äusserst genau durch mehrere Sternenpaare bestimmte. Die Sterne, mit welchen er ihn verglich, sind folgende:

\begin{center}\fontsize{10}{12}\selectfont\setlength{\tabcolsep}{5pt}\renewcommand{\arraystretch}{1.15}\begin{tabular}{c|rr|rr}
& \multicolumn{2}{c|}{Gerade Aufst.} & \multicolumn{2}{c}{Südl. Abweich.}\\\hline
I & $324^\circ$ & $34'$ & $24^\circ$ & $31'$\\
II & 323 & 10 & 24 & 59\\
III & 322 & 37 & 25 & 16\\
IV & 322 & 45 & 25 & 57\\
V & 322 & 17 & 26 & 3\\
\end{tabular}\end{center}

\Needspace{15\baselineskip}\par\indent Die Stellungen des Cometen gegen diese Sterne ergeben sich nach Hrn. Prof. \textsc{Bessel}'s Reduction, wie folgt:

\begin{center}\fontsize{9}{11}\selectfont\setlength{\tabcolsep}{3pt}\renewcommand{\arraystretch}{1.15}\begin{tabular}{l|c|c|c|rr|rr}
1812 & \shortstack{Mittlere Zeit in\\Neu-Tscherkask} & \shortstack{Vergl.\\Stern} & \shortstack{Anz. d.\\Beob.} & \multicolumn{2}{c|}{\shortstack{Gerade\\Aufsteigung}} & \multicolumn{2}{c}{Abweichung}\\\hline
Aug. 8 & $12^{\mathrm u}\;18^{\mathrm m}\;51^{\mathrm s}4$ & I & 4 & $-5'$ & $30''5$ & $-1'$ & $19''3$\\
11 & $12\;45\;21.7$ & II & 2 & $+25$ & 23.7 & $-3$ & 52.2\\
12 & $11\;42\;25.3$ & II & 7 & $+9$ & 7.5 & $-51$ & 15.3\\
12 & $12\;41\;41.0$ & III & 3 & $+41$ & 10.0 & $+1$ & 15.8\\
15 & $12\;10\;7.9$ & IV & 3 & $-18$ & 9.0 & $+11$ & 50.0\\
15 & $12\;30\;18.7$ & V & 5 & $+8$ & 39.8 & $+13$ & 0.0\\
17 & $13\;34\;8.4$ & V & 6 & $-25$ & 36.7 & $-1$ & 1.7\\
\end{tabular}\end{center}

\Needspace{12\baselineskip}\par\indent Von den Sternen selbst kommen nur drei (II, IV, V) in der \emph{Histoire céleste} vor. Herr Prof. \textsc{Bessel} reducirte die Angabe derselben, nach seinen eigenen Reductionselementen auf das Jahr 1812 und fand:

\begin{center}\fontsize{10}{12}\selectfont\setlength{\tabcolsep}{4pt}\renewcommand{\arraystretch}{1.15}\begin{tabular}{c|rrr|rrr}
& \multicolumn{3}{c|}{Gerade Aufsteig.} & \multicolumn{3}{c}{Südl. Abweich.}\\\hline
II & $323^\circ$ & $9'$ & $1''5$ & $24^\circ$ & $59'$ & $26''5$\\
IV & 322 & 44 & 32.3 & 25 & 56 & 53.4\\
V & 322 & 16 & 17.1 & 26 & 3 & 36.9\\
\end{tabular}\end{center}

\Needspace{7\baselineskip}\par\indent Die Rectascensionen von II und IV bestimmte er überdies auch noch aus eignen Beobachtungen:

\begin{center}\fontsize{10}{12}\selectfont\setlength{\tabcolsep}{4pt}\renewcommand{\arraystretch}{1.15}\begin{tabular}{lcrl}
II & $.\;.\;.\;.\;.\;.\;.\;.$ & $323^\circ\;9'\;5''3$ & (1 Beob.)\\
IV & $.\;.\;.\;.\;.\;.\;.\;.$ & $322\;44\;56.1$ & (2 Beob.)\\
\end{tabular}\end{center}

\Needspace{13\baselineskip}\par\indent Indem diese vorzugsweise zum Grunde gelegt wurden, ergaben sich also folgende beobachtete Oerter des Cometen:

\begin{center}\fontsize{9}{11}\selectfont\setlength{\tabcolsep}{3pt}\renewcommand{\arraystretch}{1.15}\begin{tabular}{l|c|rrr|rrr}
& Mittlere Zeit & \multicolumn{3}{c|}{Gerade Aufsteig.} & \multicolumn{3}{c}{Südl. Abweich.}\\\hline
Aug. 11 & $12^{\mathrm u}\;45^{\mathrm m}\;21^{\mathrm s}7$ & $323^\circ$ & $35'$ & $9''9$: & $25^\circ$ & $3'$ & $0''4$\\
12 & $11\;42\;25.3$ & 323 & 18 & 53.7 & 25 & 14 & 59.5\\
15 & $12\;10\;7.9$ & 322 & 27 & 7.9 & 25 & 44 & 45.3\\
15 & $12\;30\;18.7$ & 322 & 25 & 37.9 & 25 & 45 & 19.0\\
17 & $13\;34\;8.4$ & 321 & 51 & 21.4 & 26 & 4 & 20.7\\
\end{tabular}\end{center}

\Needspace{13\baselineskip}\par\indent Die Vergleichung dieser Beobachtungen mit den rein elliptischen Elementen gab Hrn. Prof. \textsc{Bessel} folgende Differenzen
\SourceLeaf{380}\begin{center}\fontsize{10}{12}\selectfont\setlength{\tabcolsep}{8pt}\renewcommand{\arraystretch}{1.15}\begin{tabular}{l|c|c}
& \shortstack{Gerade\\Aufsteig.} & Abweichung\\\hline
Aug. 11 & $-31''2$ & $-10''0$:\\
12 & $-55.9$ & $+113.9$\\
15 & $-86.4$ & $+76.1$\\
15 & $-10.9$ & $+101.6$\\
17 & $-51.2$ & $+57.8$\\
\end{tabular}\end{center}

\par\indent Diese Uebereinstimmung, nach einer Zwischenzeit von mehr als drei Vierteljahren, beweist, wie gut Hr. \textsc{Bessel} die Beobachtungen zu benutzen gewusst hat, unter welchen besonders diejenigen, welche er selbst mit einem vortrefflichen Heliometer angestellt hatte, von ausgezeichneter Genauigkeit waren. Man sieht überdies sogar, dass Hr. \textsc{Bessel} selbst eine noch etwas bessere Uebereinstimmung erhalten haben würde, wenn er seine Elemente noch einmal nach den spätetsen 
damaligen Beobachtungen auszufeilen der Mühe schon damals werth gehalten hätte, da diese schon etwas in demselben Sinn von den Elementen abzuweichen anfingen. Man sieht zugleich, dass die Störungen nur einen sehr kleinen Einfluss auf die Bewegung des Cometen geäussert haben können, mit dessen Berechnung Hr. \textsc{Bessel} gegenwärtig beschäftigt ist.
